\documentclass[a4paper,12pt]{article}
\usepackage[utf8]{inputenc}
\usepackage[T1]{fontenc}
\usepackage[french]{babel}
\usepackage{graphicx}
\usepackage{hyperref}
\usepackage{geometry}
\usepackage{tikz}
\usetikzlibrary{shapes.geometric, arrows, positioning, fit, calc}
\usepackage{float}

\geometry{hmargin=2.5cm,vmargin=2.5cm}

\title{\textbf{Feuille de Réponses - Projet Base de Données}\\Smart City Neo-Sousse 2030}
\author{Étudiant: [Votre Nom] \\ Module: BDD Relationnelles}
\date{\today}

\begin{document}

\maketitle

\section*{A. Partie Théorique}

\subsection*{1. Modélisation Conceptuelle (Diagramme Entité-Association)}

\begin{figure}[H]
\centering
\begin{tikzpicture}[
    scale=0.7,
    transform shape,
    entity/.style={rectangle, draw=black, thick, minimum width=2cm, minimum height=0.8cm, align=center, font=\small},
    relation/.style={diamond, draw=black, thick, aspect=2, minimum width=1.2cm, align=center, font=\small},
    attribute/.style={ellipse, draw=black, minimum width=1.2cm, align=center, font=\small},
    link/.style={-, thick}
]
% Entités
\node[entity] (arrond) {ARRONDISSEMENT};
\node[entity] (capteur) [below=4cm of arrond] {CAPTEUR};
\node[entity] (proprio) [left=4cm of capteur] {PROPRIETAIRE};
\node[entity] (mesure) [below=4cm of capteur] {MESURE};
\node[entity] (interv) [right=4cm of capteur] {INTERVENTION};
\node[entity] (tech) [right=4cm of interv] {TECHNICIEN};
% Relations
\node[relation] (loc) at ($(arrond)!0.5!(capteur)$) {Localise};
\node[relation] (possede) at ($(proprio)!0.5!(capteur)$) {Possede};
\node[relation] (genere) at ($(capteur)!0.5!(mesure)$) {Genere};
\node[relation] (subit) at ($(capteur)!0.5!(interv)$) {Subit};
\node[relation] (effectue) at ($(interv)!0.5!(tech)$) {Effectue};
% Liens & Cardinalités
\draw[link] (arrond) -- node[right, font=\footnotesize] {1,n} (loc);
\draw[link] (loc) -- node[right, font=\footnotesize] {1,1} (capteur);
\draw[link] (proprio) -- node[above, font=\footnotesize] {1,n} (possede);
\draw[link] (possede) -- node[above, font=\footnotesize] {1,1} (capteur);
\draw[link] (capteur) -- node[right, font=\footnotesize] {1,n} (genere);
\draw[link] (genere) -- node[right, font=\footnotesize] {1,1} (mesure);
\draw[link] (capteur) -- node[above, font=\footnotesize] {1,n} (subit);
\draw[link] (subit) -- node[above, font=\footnotesize] {1,1} (interv);
\draw[link] (interv) -- node[above, font=\footnotesize] {1,1} (effectue);
\draw[link] (effectue) -- node[above, font=\footnotesize] {1,n} (tech);
\end{tikzpicture}

\end{figure}

\subsection*{2. Cardinalités et Contraintes}
\begin{itemize}
    \item \textbf{Propriétaire - Capteur (1,n)} : Un propriétaire peut posséder plusieurs capteurs, mais un capteur appartient à un seul propriétaire unique.
    \item \textbf{Arrondissement - Capteur (1,n)} : Un capteur est installé dans un seul arrondissement précis. Un arrondissement contient de nombreux capteurs.
    \item \textbf{Intervention - Technicien} : 
    \begin{itemize}
        \item Cardinalité (1,n) vers Technicien pour l'intervenant.
        \item Cardinalité (1,n) vers Technicien pour le validateur.
        \item \textbf{Contrainte d'Intégrité} : Le technicien intervenant doit être strictement différent du technicien validateur ($id\_intervenant \neq id\_validateur$).
    \end{itemize}
\end{itemize}

\subsection*{3. Modèle Relationnel Normalisé}
Le schéma relationnel dérivé du diagramme E/A est le suivant :

\begin{itemize}
    \item \textbf{ARRONDISSEMENT} (\underline{id}, nom, code\_postal, superficie\_km2, population)
    \item \textbf{PROPRIETAIRE} (\underline{id}, nom, type, contact\_email, contact\_tel)
    \item \textbf{VEHICULE} (\underline{id}, plaque, type, energie, marque, modele, capacite)
    \item \textbf{CAPTEUR} (\underline{uuid}, type, latitude, longitude, statut, date\_install, \#arrond\_id, \#proprio\_id)
    \item \textbf{INTERVENTION} (\underline{id}, \#capteur\_uuid, date\_heure, nature, duree\_minutes, cout, reduction\_co2, \#tech\_interv\_id, \#tech\_valid\_id)
\end{itemize}

\subsection*{4. Graphes des Dépendances Fonctionnelles (DF)}
Pour chaque relation, nous illustrons les dépendances fonctionnelles atomiques :

\subsubsection*{Relation CAPTEUR}
\begin{center}
\begin{tikzpicture}[
    node distance=1.5cm,
    attribute/.style={rectangle, draw=black, rounded corners, minimum height=0.6cm, align=center, font=\small},
    arrow/.style={->, >=stealth, thick}
]
    \node[attribute] (uuid) {\underline{uuid}};
    \node[attribute] (attr) [right=2cm of uuid] {type, latitude, longitude, statut, date\_install};
    \node[attribute] (fks) [below=1cm of attr] {arrond\_id, proprio\_id};

    \draw[arrow] (uuid) -- (attr);
    \draw[arrow] (uuid) |- (fks);
\end{tikzpicture}
\end{center}

\subsubsection*{Relation VEHICULE}
\begin{center}
\begin{tikzpicture}[
    node distance=1.5cm,
    attribute/.style={rectangle, draw=black, rounded corners, minimum height=0.6cm, align=center, font=\small},
    arrow/.style={->, >=stealth, thick}
]
    \node[attribute] (id) {\underline{id}};
    \node[attribute] (plaque) [right=2cm of id] {plaque};
    \node[attribute] (attr) [below=1cm of plaque] {type, energie, marque, modele};

    \draw[arrow] (id) -- (plaque);
    \draw[arrow] (id) |- (attr);
    \draw[arrow] (plaque) -- (attr) node[midway, right] {DF Transitive (Si plaque est clé candidate)};
\end{tikzpicture}
\end{center}

\subsubsection*{Relation INTERVENTION}
\begin{center}
\begin{tikzpicture}[
    node distance=1.5cm,
    attribute/.style={rectangle, draw=black, rounded corners, minimum height=0.6cm, align=center, font=\small},
    arrow/.style={->, >=stealth, thick}
]
    \node[attribute] (id) {\underline{id}};
    \node[attribute] (details) [right=2cm of id] {date, nature, duree, cout};
    \node[attribute] (refs) [below=1cm of details] {capteur\_uuid, tech\_intervenant, tech\_validateur};

    \draw[arrow] (id) -- (details);
    \draw[arrow] (id) |- (refs);
\end{tikzpicture}
\end{center}

\subsection*{4. Analyse des Redondances et Normalisation}
\begin{itemize}
    \item \textbf{1FN}: Tous les attributs sont atomiques.
    \item \textbf{2FN}: Toutes les tables ont une clé primaire simple, ou dépendent de toute la clé.
    \item \textbf{3FN}: Pas de dépendance transitive (ex: adresse proprio dans table Proprietaire).
    \item Le schéma est en 3FN/BCNF.
\end{itemize}

\end{document}
